\documentclass[10pt,a4paper]{amsart}
\usepackage[margin=19mm]{geometry}
\usepackage{amsmath,amssymb,mathtools}
\usepackage[scr=rsfs,cal=euler]{mathalfa}
\usepackage{xfrac}
\usepackage[unicode,hidelinks,bookmarks=false]{hyperref}
\newcommand{\ie}{i.e.\ }
\newcommand{\cf}{cf.\ }
\newcommand{\resp}{respectively\ }
\newcommand{\Subsection}{\subsection}
\providecommand{\nobreakdash}{\nobreak\hbox{-}}
\setlength{\emergencystretch}{2em}
\multlinegap0pt
\usepackage{bm}
\usepackage{bbm}
\usepackage{dsfont}
\usepackage{xspace}
\usepackage{cancel}
\usepackage{mathtools}
\usepackage{amsthm}
\makeatletter
\@ifundefined{thm}{\newtheorem{thm}{Theorem}}{}
\@ifundefined{lem}{\newtheorem{lem}[thm]{Lemma}}{}
\makeatother
\makeatletter
\expandafter\ifx\csname customThm\endcsname\relax
\newenvironment{customThm}[1]{\paragraph*{\textbf{Theorem #1.}}\itshape}{\par}
\fi
\makeatother
\title[A thermodynamic path metric for complex H\'enon maps]{A thermodynamic path metric for complex H\'enon maps}
\author{Fabrizio Bianchi, Yan Mary He}
\date{Source: arXiv:2606.29363v1 (CC BY 4.0)}
\begin{document}
\maketitle

\noindent\emph{Source and licence.} A thermodynamic path metric for complex H\'enon maps. Version arXiv:2606.29363v1 (CC BY 4.0), distributed under the Creative Commons Attribution 4.0 International licence, as recorded in the arXiv licence metadata for this exact version and shown on the abstract page https://arxiv.org/abs/2606.29363v1. Full rights notice: licenses/arxiv-2606.29363-CC-BY-4.0.txt.\par\smallskip

\noindent\textit{Editorial scope.} Counted unit: the Thm whose source label is \texttt{t:mityagin} in the article (source0.tex lines 1298--1305, source label \texttt{t:mityagin}), delivered together with the article's own complete original proof of it (source0.tex lines 1308--1359, one \texttt{proof} environment of 52 lines). No mathematics is written, repaired or re-derived: statement and proof are the article's own text line for line. Required context retained uncounted: lem:semi-reduction (lines 1287--1292). Every definition, display and label of those lines is retained unchanged. Omitted from this package and not redistributed: 1--1286, 1293--1297, 1306--1307, 1360--1392. Nothing inside a retained excerpt is omitted or altered: every retained body is delivered exactly as the article's lines are written, so no omission receipt of any kind accompanies this package. Citations: the retained excerpts carry no citation command at all, so no bibliography entry of the article is transcribed and no reference is redistributed. Reference adaptation: none was needed and none was made; every reference inside the delivered unit resolves inside it, so no unresolved reference or citation is produced. Printed slips kept literal: no printed word, symbol, hypothesis or number of the counted statement or of its proof was altered. Layout and class: the article's own class is \texttt{amsart} with packages including inputenc, biblatex, amssymb, amsfonts, amsmath, amscd, xy, enumerate, mathrsfs, graphicx, fontenc, color, hyperref, dsfont; the wrapper is the lane's pinned \texttt{amsart} editorial wrapper at 10pt on A4, which restates the theorem-like environments the retained text needs and adds the article's own macro definitions that the retained text needs (none), with extra wrapper packages (bm, bbm, dsfont, xspace, cancel, mathtools). The delivered numbering is the unit's own. Assets: no retained span contains a figure environment, a graphics-inclusion command, a PDF-page inclusion, a table or a TikZ picture; no image file is copied into this package, the package declares no asset dependency and no image file is redistributed with it. Licence: version arXiv:2606.29363v1 of this article is distributed under the Creative Commons Attribution 4.0 International licence, as recorded in the arXiv licence metadata for this exact version and shown on the abstract page https://arxiv.org/abs/2606.29363v1. A full rights notice is delivered at licenses/arxiv-2606.29363-CC-BY-4.0.txt. Characters: every retained excerpt is pure ASCII, so the delivered engine, XeLaTeX with its default Latin Modern text font, reports no missing character.
\begin{lem}\label{lem:semi-reduction}
Let $M$ be a real-analytic manifold and let $Q$ be a real-analytic
positive semi-definite quadratic form on $TM$. Assume that every non-constant
piecewise $C^1$-path in $M$ has positive $Q$-length. Then the path
pseudo-distance induced by $Q$ separates points locally.
\end{lem}

\begin{thm}[Mityagin]\label{t:mityagin}
Let $V\subset\mathbb R^N$ be open and let
$F\colon V\to\mathbb R$
be a real-analytic function which is not identically zero. Then the zero set
$\{F=0\}$
is covered by a countable union of 
(not necessarily closed) real-analytic submanifolds of $V$.
\end{thm}

\begin{proof}[Proof of Lemma \ref{lem:semi-reduction}]
Since the statement is local, we work in a coordinate ball
$V\subset\mathbb R^N$.
Let $B(x)$ be the positive semi-definite symmetric matrix representing $Q$
in this chart, i.e.,
$Q(x,\xi)=\xi^T B(x)\xi$.
The entries of $B$ are real-analytic.

We first claim that
$\Delta(x):=\det B(x)$
is not identically zero on any open subset of $V$. Indeed, if
$\Delta\equiv0$ on some open set $W$, then $B(x)$ has a non-trivial kernel
for every $x\in W$. On a smaller open subset $W'\subset W$, the rank of
$B(x)$ is constant. Hence $\ker B(x)$ forms a non-trivial real-analytic
subbundle of $TW'$. Choosing a non-vanishing real-analytic vector field
$X(x)\in\ker B(x)$ and integrating it gives a non-constant analytic path
$\eta$ such that
$Q(\eta(t),\dot\eta(t))=0$
for all $t$. This path has zero length, contradicting the assumption.
Therefore $\Delta$ is not identically zero.

It follows that $Q$ is positive definite on
$V\setminus Z_1$ where $Z_1:=\{\Delta=0\}$.
By Theorem
\ref{t:mityagin}, $Z_1$ is covered by countably many real-analytic
submanifolds of positive codimension. Any zero-distance degeneracy class must
be contained in $Z_1$; otherwise it would meet a region where $Q$ is
positive definite, which 
would give positive
distance.

Now restrict $Q$ to each analytic submanifold $Y\subset Z_1$. The restricted
form is again real-analytic and positive semi-definite on $TY$. If its
determinant were identically zero on an open subset of $Y$, the same
constant-rank argument would produce a non-constant path inside $Y$ of zero
length, again contradicting the assumption. Therefore the determinant of the
restricted form is not identically zero. Applying 
Theorem  \ref{t:mityagin} again, 
its
zero locus inside $Y$ is covered by countably many analytic submanifolds of
strictly smaller dimension.

By induction, we can reduce 
the problem to the case of dimension $1$, i.e., an analytic path.
In this case,
a non-zero real-analytic
non-negative quadratic form has only isolated zeros unless it vanishes
identically; the latter is impossible by the positive-length assumption.
Therefore two distinct points in the same one-dimensional analytic stratum have
positive path distance. This
proves local separation and completes the proof.
\end{proof}

\end{document}
